\documentclass[a0paper,portrait,fontscale=0.27]{baposter}
%\renewcommand{\baselinestretch}{2.2}
%fontscale=0.292
%showframe
\usepackage{setspace}
\usepackage{relsize}		% For \smaller
\usepackage{url}			% For \url
% \usepackage{epstopdf}	
\usepackage{amssymb}
% \usepackage{multicol}

%\usepackage{pstricks}
\usepackage{amsmath}
% Included EPS files automatically converted to PDF to include with pdflatex
%\usepackage{hyperref}



%%% Global Settings %%%%%%%%%%%%%%%%%%%%%%%%%%%%%%%%%%%%%%%%%%%%%%%%%%%%%%%%%%%

\graphicspath{{pix/}}	% Root directory of the pictures 
\tracingstats=2			% Enabled LaTeX logging with conditionals

%%% Color Definitions %%%%%%%%%%%%%%%%%%%%%%%%%%%%%%%%%%%%%%%%%%%%%%%%%%%%%%%%%

\definecolor{bordercol}{RGB}{34,39,66}
% \definecolor{headercol1}{RGB}{186,215,230}
% \definecolor{headercol2}{RGB}{80,80,80} %7	62	137
% \definecolor{headerfontcol}{RGB}{0,0,0}
% \definecolor{boxcolor}{RGB}{186,215,230}



\definecolor{boxcolorsimple}{RGB}{186,231,185} 
\definecolor{boxcolordetails}{RGB}{170,205,213}
\definecolor{boxcolorexplanation}{RGB}{205,190,160}
\definecolor{boxcolorexamples}{RGB}{170,205,213}
\definecolor{headercolfade}{RGB}{34,39,66}



% \definecolor{highlightcolor}{RGB}{0,0,0}

\newcommand{\hlt}[1]{{\bf\color{highlightcolor} #1}}

\newcommand{\startblockmodifiers}{\sf
	%\large
	\vspace{-0.4em}
}

% ,
% 	boxColorOne=boxcolorsimple,
% 	headerColorOne=boxcolorsimple, 
% 	headerColorTwo=headercolfade
\newcommand{\blockbasic}[3]{
	\headerbox{#1}{name=#2,#3}{
	\startblockmodifiers
	\input{#2}
	}
}


\newcommand{\blockdetail}[3]{
	\headerbox{#1}{name=#2,#3,boxColorOne=boxcolordetails,
	headerColorOne=boxcolordetails, headerColorTwo=headercolfade}{
	\startblockmodifiers
	\input{#2}
	%\endblockmodifiers
	}
}

\newcommand{\blockexplain}[3]{
	\headerbox{#1}{name=#2,#3,
	boxColorOne=boxcolorexplanation,
	headerColorOne=boxcolorexplanation, 
	headerColorTwo=headercolfade}{
	\startblockmodifiers
	\input{#2}
	%\endblockmodifiers
	}
}


\newcommand{\blockexample}[3]{
	\headerbox{#1}{name=#2,#3,
	boxColorOne=boxcolorexamples,
	headerColorOne=boxcolorexamples, 
	headerColorTwo=headercolfade}{
	\startblockmodifiers
	\input{#2}
	%\endblockmodifiers
	}
}

\newcommand{\TODO}[1]{{\bf\color{highlightcolor} TODO: #1}}

\newcommand{\fig}[3]{
%{\centering{\vspace{0.2em}
%\smaller\it
%	#1}
\begin{center}
\vspace{-0.5em}
\smaller\it
#1
\includegraphics[height=#3em]{#2}
\vspace{-0.9em}
\end{center}
}

\newcommand{\figstwo}[5]{

\begin{center}
%\smaller
\vspace{-0.5em}
\begin{tabular}{cc}
%\vspace{-0.5em}
\smaller\it
#1
& 
\smaller\it
#3
\\

\includegraphics[height=#5em]{#2}
%\vspace{-0.9em}
&
%\vspace{-0.5em}
\includegraphics[height=#5em]{#4}
%\vspace{-0.9em}
\end{tabular}
\end{center}
}

%\newcommand{\eqref}[1]{(\ref{#1})}

\newcommand{\norm}[2]{\left\| #1 \right\|_{#2}}


\newcommand{\normeddif}[4]{
\norm{p_{#1}^{#2} - p_{#3}^{#4}}{L^2}
}


%%%%%%%%%%%%%%%%%%%%%%%%%%%%%%%%%%%%%%%%%%%%%%%%%%%%%%%%%%%%%%%%%%%%%%%%%%%%%%%%
%%% Utility functions %%%%%%%%%%%%%%%%%%%%%%%%%%%%%%%%%%%%%%%%%%%%%%%%%%%%%%%%%%

%%% Save space in lists. Use this after the opening of the list %%%%%%%%%%%%%%%%
\newcommand{\compresslist}{
	\smaller
	\setlength{\itemsep}{0pt}
	\setlength{\parskip}{0pt}
	\setlength{\parsep}{0pt}
}
\renewcommand*\rmdefault{ppl}\normalfont\upshape
%\renewcommand*\rmdefault{iwona}\normalfont\upshape
%\usepackage{helvet}

%%%%%%%%%%%%%%%%%%%%%%%%%%%%%%%%%%%%%%%%%%%%%%%%%%%%%%%%%%%%%%%%%%%%%%%%%%%%%%%
%%% Document Start %%%%%%%%%%%%%%%%%%%%%%%%%%%%%%%%%%%%%%%%%%%%%%%%%%%%%%%%%%%%
%%%%%%%%%%%%%%%%%%%%%%%%%%%%%%%%%%%%%%%%%%%%%%%%%%%%%%%%%%%%%%%%%%%%%%%%%%%%%%%

\begin{document}
\typeout{Poster rendering started}

%%% Setting Background Image %%%%%%%%%%%%%%%%%%%%%%%%%%%%%%%%%%%%%%%%%%%%%%%%%%
%\background{
%	\begin{tikzpicture}[remember picture,overlay]%
%	\draw (current page.north west)+(-2em,2em) node[anchor=north west]
%	{\includegraphics[height=1.1\textheight]{background}};
%	\end{tikzpicture}
%}

%%% General Poster Settings %%%%%%%%%%%%%%%%%%%%%%%%%%%%%%%%%%%%%%%%%%%%%%%%%%%
%%%%%% Eye Catcher, Title, Authors and University Images %%%%%%%%%%%%%%%%%%%%%%
\begin{poster}{
	grid=false,
	% Option is left on true though the eyecatcher is not used. The reason is
	% that we have a bit nicer looking title and author formatting in the headercol
	% this way
	eyecatcher=true, 
	borderColor=bordercol,
% 	headerColorOne=headercol1,
% 	headerColorTwo=headercol2,
% 	headerFontColor=headerfontcol,
	% Only simple background color used, no shading, so boxColorTwo isn't necessary
	boxColorOne=bordercol,
% 	headershape=roundedright,
	headerfont=\Large\bf, 
	textborder=rectangle,
	background=none,
	headerborder=open,
    boxshade=plain
	%textfont=\sf
}
%\rm Roman.
%\it Italics.
%\em Emphasis (toggles between \it and \rm).
%\emph (LaTeX2e emphasis command)
%\bf Boldface.
%\sl Slanted.
%\sf Sans serif.
%\sc Small caps.
%\tt Typewriter.
%%% Eye Cacther %%%%%%%%%%%%%%%%%%%%%%%%%%%%%%%%%%%%%%%%%%%%%%%%%%%%%%%%%%%%%%%
{
	%Eye Catcher, empty if option eyecatcher=false - unused
%\setlength\fboxsep{0pt}
%\setlength\fboxrule{0.5pt}
%	\fbox{
%		\begin{minipage}{14em}
%			\includegraphics[height=8em]{teaserCrop.pdf}
			%Replace me
%		\end{minipage}
%	}

}
%%% Title %%%%%%%%%%%%%%%%%%%%%%%%%%%%%%%%%%%%%%%%%%%%%%%%%%%%%%%%%%%%%%%%%%%%%
{\bf
	%\hspace{1.7em}
%\hspace{+0.9em}
	%\renewcommand{\baselinestretch}{2.2}
	%\raggedright
	%\textls
    % Sequential multi-realization probabilistic interpretation of well logs by a deep-learning method
    % Probabilistic log interpretation and extrapolation with deep learning
    AI-based multi-modal interpretation of logs
    % with AI \\
    for ahead-of-bit probabilistic ROP prediction 
}
%%% Authors %%%%%%%%%%%%%%%%%%%%%%%%%%%%%%%%%%%%%%%%%%%%%%%%%%%%%%%%%%%%%%%%%%%
%\rm Roman.
%\it Italics.
%\em Emphasis (toggles between \it and \rm).
%\emph (LaTeX2e emphasis command)
%\bf Boldface.
%\sl Slanted.
%\sf Sans serif.
%\sc Small caps.
%\tt Typewriter.
{\sf
	\vspace{0.8em} 
	Adrian Ambrus, Sergey Alyaev, Nazanin Jahani, Ahmed H. Elsheikh\\
		{\smaller\tt \setlength{\unitlength}{1em} 
		E-mail:aamb@norceresearch.no}
}
%%% Logo %%%%%%%%%%%%%%%%%%%%%%%%%%%%%%%%%%%%%%%%%%%%%%%%%%%%%%%%%%%%%%%%%%%%%%
{
% \setlength\fboxsep{0pt}
% \setlength\fboxrule{0.5pt}
% 	\fbox{
% 		\begin{minipage}{8em}
			\centering
			\includegraphics[height=8em]{digiwells_logo_negative.png}
% 			\\
% 			\includegraphics[height=2em]{research-council-logo.png}
% 		\end{minipage}
% 	}
}


%%% The poster %%%%%%%%%%%%%%%%%%%%%%%%%%%%%%%%%%%%%%%%%%%%%%%%%%%%%%%%%%%%%%%%%
%The boxes start here

\blockbasic{Motivation}{motivation}{column=0,row=0}

%\blockbasic{Model problem}{modelProblem}{column=0,below=motivation}

\blockexplain{Method}{method}{column=1,row=0,span=2}

%end of column 0


\blockdetail{Dataset description}{datasetDescription}{column=0,below=motivation}

\blockexplain{Training details}{trainingDetails}{column=0,below=datasetDescription,above=bottom}

%%%%%%%%%%%%%%%%%%%%%%%%%%%%

\blockexample{Numerical results}{numericalResults}{column=1,span=2,below=method}

%\blockexample{Numerical results}{numericalResults2}{column=2,below=errorContributions}


\blockbasic{Conclusions}{conclusion}{column=1,span=2,below=numericalResults}


\blockexplain{References}{references}{column=1,span=2,below=conclusion,above=bottom}

%\blockexplain{ROP prediction}{ropPrediction}{column=2,row=0,above=numericalResults}


%\blockbasic{Aknowledgements}{aknowledgements}{column=2,below=references,above=numericalResults}



%%% The end of the poster %%%%%%%%%%%%%%%%%%%%%%%%%%%%%%%%%%%%%%%%%%%%%%%%%%%%%%


%\headerbox{Basic Concepts}{name=definitions,column=0,below=motivation}{
%\begin{equation}
%	123
%	\label{eq123}
%\end{equation}
%}

%\headerbox{Models}{name=models,column=0,below=definitions}{
%}

%\headerbox{References}{name=references,column=0,below=models}{
%%\smaller													% Make the whole text smaller
%%\vspace{-0.4em} 										% Save some space at the beginning
%%\bibliographystyle{plain}							% Use plain style
%%\renewcommand{\section}[2]{\vskip 0.05em}		% Omit "References" title
%%\begin{thebibliography}{1}							% Simple bibliography with widest label of 1
%%\itemsep=-0.01em										% Save space between the separation
%%\setlength{\baselineskip}{0.4em}					% Save space with longer lines
%%\bibitem{prevWork1} Laszlo Gulyas, Richard Legendi: \emph{Effects of Sample Duration on Network Statistics in Elementary Models of Dynamic Networks}, International Conference on Computational Science, Singapore (2011) 
%%\bibitem{prevWork2} Laszlo Gulyas, Susan Khor, Richard Legendi and George Kampis \emph{Cumulative Properties of Elementary Dynamic Networks}, The International Sunbelt Social Network Conference XXXI (2011)
%%\bibitem{gulya-kampis1} Gulyas, Laszlo et al.: \emph{Betweenness Centrality Dynamics in Networks of Changing Density}. Presented at the 19th International Symposium on Mathematical Theory of Networks and Systems (MTNS 2010)
%%\end{thebibliography}
%}

%\headerbox{Acknowledgements}{name=acknowledgements,column=0,below=references, above=bottom,boxColorOne=red, boxColorTwo=black}{
%\smaller						% Make the whole text smaller
%\vspace{-0.4em}			% Save some space at the beginning
%} 

%\headerbox{Dynamic Networks are Sensitive to Aggregation}{name=density,span=2,column=1,row=0,headerColorOne=black,headerColorTwo=white}{
%}

%%\headerbox{Degree Distribution Radically Changes}
%%{name=degreeDistribution,span=2,column=1,below=density,above=bottom}{
%%}

%\blockbasic{End}{introduction}{span=2,column=1,below=density,above=bottom}

\end{poster}
\end{document}
